\documentclass[11pt,a4paper]{article}

\usepackage[margin=2.2cm]{geometry}
\usepackage{kotex}
\usepackage{amsmath,amssymb}
\usepackage{booktabs}
\usepackage{array}
\usepackage{xcolor}
\usepackage{listings}
\usepackage[most]{tcolorbox}
\usepackage{enumitem}
\usepackage{hyperref}
\usepackage{fancyhdr}
\usepackage{titlesec}

\definecolor{codebg}{RGB}{247,247,247}
\definecolor{codeframe}{RGB}{210,210,210}
\definecolor{keywordcolor}{RGB}{0,70,160}
\definecolor{commentcolor}{RGB}{0,120,80}
\definecolor{stringcolor}{RGB}{170,50,50}
\definecolor{titleblue}{RGB}{35,75,130}

\hypersetup{colorlinks=true,linkcolor=titleblue,urlcolor=titleblue}

\pagestyle{fancy}
\fancyhf{}
\lhead{C++ Programming}
\rhead{Day 7}
\cfoot{\thepage}

\titleformat{\section}{\Large\bfseries\color{titleblue}}{\thesection.}{0.6em}{}
\titleformat{\subsection}{\large\bfseries}{\thesubsection.}{0.5em}{}

\lstdefinestyle{cppstyle}{
    language=C++,
    backgroundcolor=\color{codebg},
    frame=single,
    rulecolor=\color{codeframe},
    basicstyle=\ttfamily\small,
    keywordstyle=\color{keywordcolor}\bfseries,
    commentstyle=\color{commentcolor},
    stringstyle=\color{stringcolor},
    numbers=left,
    numberstyle=\tiny\color{gray},
    numbersep=8pt,
    tabsize=4,
    showstringspaces=false,
    breaklines=true,
    columns=fullflexible
}
\lstset{style=cppstyle}

\newtcolorbox{learningbox}{colback=blue!4,colframe=titleblue,title=\textbf{학습 목표},fonttitle=\bfseries}
\newtcolorbox{importantbox}{colback=yellow!8,colframe=orange!70!black,title=\textbf{핵심 포인트},fonttitle=\bfseries}
\newtcolorbox{practicebox}{colback=red!3,colframe=red!55!black,title=\textbf{실습},fonttitle=\bfseries}

\begin{document}

\begin{titlepage}
    \centering
    \vspace*{3cm}
    {\Huge\bfseries C++ Programming\par}
    \vspace{0.8cm}
    {\LARGE Day 7: 객체 복사와 자원 관리\par}
    \vspace{2cm}
    {\large 복사 생성자, 얕은 복사, 깊은 복사, 복사 대입, Rule of Three, RAII\par}
    \vfill
\end{titlepage}

\tableofcontents
\newpage

\section{학습 목표}
\begin{learningbox}
이 장을 마치면 다음 내용을 설명하고 직접 코드로 작성할 수 있다.
\begin{itemize}[leftmargin=1.5em]
    \item 복사 생성자가 호출되는 상황 구분하기
    \item 컴파일러가 만드는 기본 복사 생성자의 동작 이해하기
    \item 얕은 복사에서 포인터 주소가 공유되는 이유 설명하기
    \item 깊은 복사를 사용하여 독립적인 힙 메모리 만들기
    \item 복사 생성자와 복사 대입 연산자의 차이 구분하기
    \item 자기 대입 문제를 검사하고 안전하게 처리하기
    \item Rule of Three의 세 함수를 함께 구현하기
    \item RAII를 사용하여 객체 수명과 자원 수명 연결하기
    \item \texttt{private} 멤버와 \texttt{const} 멤버 함수 사용하기
\end{itemize}
\end{learningbox}

\section{복사 생성자}
복사 생성자(copy constructor)는 이미 존재하는 객체를 이용하여 새로운 객체를 만들 때 호출되는 특별한 생성자이다. 일반적인 형태는 다음과 같다.

\begin{lstlisting}
ClassName(const ClassName& other)
{
}
\end{lstlisting}

예를 들어 \texttt{Book} 객체 \texttt{b1}을 이용하여 새로운 객체 \texttt{b2}를 만들면 복사 생성자가 호출된다.

\begin{lstlisting}
Book b1("C++");
Book b2 = b1;
\end{lstlisting}

다음 표현도 같은 의미이다.

\begin{lstlisting}
Book b2(b1);
\end{lstlisting}

\subsection{왜 참조를 사용하는가}
복사 생성자의 매개변수는 일반적으로 다음과 같이 참조로 받는다.

\begin{lstlisting}
Book(const Book& other)
\end{lstlisting}

만약 값으로 받는다면 매개변수 \texttt{other}를 만들기 위해 다시 복사 생성자가 필요하다. 그 복사 생성자를 호출하려면 또 다른 복사가 필요하므로 무한 재귀가 발생한다.

\begin{lstlisting}
Book(Book other)
{
}
\end{lstlisting}

따라서 복사 생성자는 원본 객체를 참조로 받아야 한다.

\subsection{왜 \texttt{const}를 사용하는가}
\texttt{const}는 복사 과정에서 원본 객체를 변경하지 않겠다는 뜻이다. 또한 \texttt{const} 객체도 복사할 수 있게 한다.

\begin{lstlisting}
const Book b1("C++");
Book b2 = b1;
\end{lstlisting}

\begin{importantbox}
복사 생성자의 표준 형태는 \texttt{ClassName(const ClassName\& other)}이다. 참조는 불필요한 복사를 막고, \texttt{const}는 원본 객체를 보호한다.
\end{importantbox}

\subsection{전체 예제}
\begin{lstlisting}
#include <iostream>
#include <string>

using namespace std;

class Book
{
public:
    string title;

    Book(string title)
    {
        this->title = title;
    }

    Book(const Book& other)
    {
        cout << "Book Copied" << endl;
        this->title = other.title;
    }

    void showInfo() const
    {
        cout << "Title: " << this->title << endl;
    }
};

int main()
{
    Book b1("C++");
    Book b2 = b1;

    b1.title = "Python";

    b1.showInfo();
    b2.showInfo();

    return 0;
}
\end{lstlisting}

\texttt{std::string}은 내부적으로 안전한 복사를 제공한다. 따라서 \texttt{b1.title}을 변경해도 \texttt{b2.title}은 함께 바뀌지 않는다.

\section{얕은 복사}
얕은 복사(shallow copy)는 객체가 가진 멤버 값을 그대로 복사한다. 일반적인 정수나 문자열에서는 문제가 드러나지 않을 수 있지만, 포인터 멤버에서는 주소 자체가 복사된다.

\subsection{동적 메모리를 가진 클래스}
\begin{lstlisting}
class Student
{
public:
    int* score;

    Student(int value)
    {
        score = new int(value);
    }

    ~Student()
    {
        delete score;
    }
};
\end{lstlisting}

\texttt{score}는 정수 값 자체가 아니라 힙 메모리의 주소를 저장한다.

\begin{verbatim}
score ----> [100]
\end{verbatim}

\subsection{기본 복사 생성자의 동작}
사용자가 복사 생성자를 작성하지 않으면 컴파일러가 기본 복사 생성자를 만든다. 개념적으로 다음과 비슷하게 동작한다.

\begin{lstlisting}
Student(const Student& other)
{
    score = other.score;
}
\end{lstlisting}

이 코드는 새로운 정수 메모리를 만들지 않는다. 단지 원본 포인터에 저장된 주소를 그대로 복사한다.

\begin{verbatim}
s1.score ----\
              ---> [100]
s2.score ----/
\end{verbatim}

두 객체는 하나의 힙 메모리를 공유한다.

\subsection{값 변경 문제}
\begin{lstlisting}
Student s1(100);
Student s2 = s1;

*s2.score = 999;
\end{lstlisting}

\texttt{s1.score}와 \texttt{s2.score}가 같은 메모리를 가리키므로 두 객체에서 모두 \texttt{999}가 출력된다.

\subsection{이중 해제 문제}
프로그램이 끝나면 두 객체의 소멸자가 각각 호출된다.

\begin{lstlisting}
delete s2.score;
delete s1.score;
\end{lstlisting}

하지만 두 포인터가 같은 메모리를 가리키므로 같은 메모리를 두 번 해제하게 된다. 이것을 이중 해제(double delete)라고 하며, 프로그램의 동작은 정의되지 않는다.

\begin{importantbox}
포인터 멤버를 기본 복사하면 포인터가 가리키는 값이 아니라 주소가 복사된다. 따라서 값 공유, 댕글링 포인터, 이중 해제 문제가 발생할 수 있다.
\end{importantbox}

\subsection{얕은 복사 전체 예제}
\begin{lstlisting}
#include <iostream>

using namespace std;

class Student
{
public:
    int* score;

    Student(int value)
    {
        score = new int(value);
    }

    ~Student()
    {
        delete score;
    }

    void showInfo() const
    {
        cout << *score << endl;
    }
};

int main()
{
    Student s1(100);
    Student s2 = s1;

    *s2.score = 999;

    s1.showInfo();
    s2.showInfo();

    return 0;
}
\end{lstlisting}

이 예제는 얕은 복사의 문제를 보여 주기 위한 코드이다. 실제 프로그램에서는 그대로 사용하면 안 된다.

\section{깊은 복사}
깊은 복사(deep copy)는 포인터의 주소만 복사하지 않고 새로운 힙 메모리를 만든 뒤 원본 값을 복사한다.

\begin{lstlisting}
Book(const Book& other)
{
    price = new int(*other.price);
}
\end{lstlisting}

\texttt{other.price}는 원본 객체의 포인터이고, \texttt{*other.price}는 그 포인터가 가리키는 실제 정수 값이다.

\begin{lstlisting}
new int(*other.price)
\end{lstlisting}

위 코드는 원본 값으로 초기화된 새로운 정수 공간을 힙에 만든다.

\begin{verbatim}
b1.price ----> [30000]
b2.price ----> [30000]
\end{verbatim}

두 값은 같지만 메모리 주소는 서로 다르다. 따라서 한 객체의 값을 변경해도 다른 객체에는 영향을 주지 않는다.

\subsection{전체 예제}
\begin{lstlisting}
#include <iostream>

using namespace std;

class Book
{
public:
    int* price;

    Book(int price)
    {
        this->price = new int(price);
    }

    Book(const Book& other)
    {
        cout << "Book Deep Copied" << endl;
        this->price = new int(*other.price);
    }

    ~Book()
    {
        delete this->price;
    }

    void showInfo() const
    {
        cout << *this->price << endl;
    }
};

int main()
{
    Book b1(30000);
    Book b2 = b1;

    *b2.price = 50000;

    b1.showInfo();
    b2.showInfo();

    return 0;
}
\end{lstlisting}

출력 결과는 다음과 같다.

\begin{verbatim}
Book Deep Copied
30000
50000
\end{verbatim}

\section{복사 대입 연산자}
복사 생성자는 새 객체를 만들 때 사용된다. 이미 존재하는 객체에 다른 객체를 대입할 때는 복사 대입 연산자(copy assignment operator)가 호출된다.

\subsection{복사 생성과 복사 대입 구분}
\begin{lstlisting}
Book b1(30000);
Book b2 = b1;   // Copy constructor

Book b3(50000);
b3 = b1;        // Copy assignment operator
\end{lstlisting}

첫 번째 경우에는 \texttt{b2}가 새로 생성된다. 두 번째 경우에는 \texttt{b3}가 이미 존재하므로 기존 자원을 정리하고 새로운 값을 복사해야 한다.

\subsection{기본 형태}
\begin{lstlisting}
Book& operator=(const Book& other)
{
    return *this;
}
\end{lstlisting}

반환형이 \texttt{Book\&}인 이유는 현재 객체 자신을 참조로 반환하여 연속 대입을 가능하게 하기 위해서이다.

\begin{lstlisting}
b3 = b2 = b1;
\end{lstlisting}

\subsection{기존 메모리를 정리해야 하는 이유}
이미 존재하는 객체는 자신이 소유한 힙 메모리를 가지고 있다. 새 주소를 바로 덮어쓰면 기존 주소를 잃어버려 메모리 누수가 발생한다.

잘못된 형태는 다음과 같다.

\begin{lstlisting}
this->price = new int(*other.price);
\end{lstlisting}

기존 메모리를 정리한 뒤 새 메모리를 연결해야 한다. 더 안전한 순서는 먼저 새 메모리를 확보하고 그 다음 기존 메모리를 해제하는 것이다.

\begin{lstlisting}
int* newPrice = new int(*other.price);

delete this->price;
this->price = newPrice;
\end{lstlisting}

새 메모리 할당에 실패하더라도 기존 데이터가 먼저 사라지지 않는다는 장점이 있다.

\subsection{자기 대입 문제}
다음 코드에서는 원본 객체와 대입받는 객체가 같다.

\begin{lstlisting}
b1 = b1;
\end{lstlisting}

이때 \texttt{this}와 \texttt{\&other}는 같은 주소를 가진다. 기존 메모리를 먼저 삭제한 뒤 \texttt{other.price}를 읽으면 이미 삭제된 메모리를 역참조하게 된다. 따라서 자기 대입 여부를 먼저 검사한다.

\begin{lstlisting}
if (this == &other)
{
    return *this;
}
\end{lstlisting}

\subsection{전체 예제}
\begin{lstlisting}
#include <iostream>

using namespace std;

class Book
{
public:
    int* price;

    Book(int price)
    {
        this->price = new int(price);
    }

    Book(const Book& other)
    {
        this->price = new int(*other.price);
    }

    Book& operator=(const Book& other)
    {
        if (this == &other)
        {
            return *this;
        }

        int* newPrice = new int(*other.price);

        delete this->price;
        this->price = newPrice;

        cout << "Copy Assignment" << endl;

        return *this;
    }

    ~Book()
    {
        delete this->price;
    }

    void showInfo() const
    {
        cout << *this->price << endl;
    }
};

int main()
{
    Book b1(30000);
    Book b2(50000);

    b2 = b1;

    *b2.price = 100000;

    b1.showInfo();
    b2.showInfo();

    return 0;
}
\end{lstlisting}

출력 결과는 다음과 같다.

\begin{verbatim}
Copy Assignment
30000
100000
\end{verbatim}

\section{Rule of Three와 RAII}
동적 메모리처럼 직접 관리해야 하는 자원을 가진 클래스는 복사와 소멸 동작을 함께 설계해야 한다.

\subsection{Rule of Three}
다음 세 함수 중 하나를 직접 구현해야 한다면 나머지 둘도 필요할 가능성이 높다.

\begin{enumerate}[leftmargin=2em]
    \item 소멸자
    \item 복사 생성자
    \item 복사 대입 연산자
\end{enumerate}

이를 Rule of Three라고 한다.

\begin{lstlisting}
~ClassName();
ClassName(const ClassName& other);
ClassName& operator=(const ClassName& other);
\end{lstlisting}

소멸자만 구현하면 복사 과정에서 주소가 공유될 수 있다. 복사 생성자만 구현하면 이미 존재하는 객체끼리의 대입에서 얕은 복사가 발생할 수 있다. 복사 대입 연산자만 구현하면 새 객체를 복사 생성할 때 기본 복사가 사용될 수 있다. 따라서 세 함수를 하나의 세트로 생각해야 한다.

\subsection{RAII}
RAII는 Resource Acquisition Is Initialization의 약자이다. 핵심은 객체가 생성될 때 자원을 획득하고 객체가 소멸될 때 자원을 자동으로 해제하도록 만드는 것이다.

\begin{lstlisting}
NumberBox(int value)
{
    number = new int(value);
}

~NumberBox()
{
    delete number;
}
\end{lstlisting}

객체의 수명과 힙 메모리의 수명이 연결된다.

\begin{verbatim}
객체 생성 -> 자원 획득
객체 사용 -> 자원 유지
객체 소멸 -> 자원 해제
\end{verbatim}

RAII는 동적 메모리뿐 아니라 파일, 네트워크 연결, 락과 같은 자원 관리에도 사용된다.

\subsection{왜 멤버를 \texttt{private}으로 두는가}
포인터 멤버가 \texttt{public}이면 외부 코드가 객체의 자원 관리 규칙을 깨뜨릴 수 있다.

\begin{lstlisting}
delete box.number;
box.number = nullptr;
\end{lstlisting}

이런 조작은 이중 해제나 메모리 누수를 만들 수 있다. 따라서 자원 포인터는 \texttt{private}으로 숨기고, 필요한 기능만 \texttt{public} 멤버 함수로 제공한다.

\begin{lstlisting}
private:
    int* number;

public:
    void setNumber(int value);
    void showInfo() const;
\end{lstlisting}

이것은 데이터와 구현 세부 사항을 숨기는 캡슐화의 한 예이다.

\subsection{\texttt{const} 멤버 함수}
멤버 함수 뒤의 \texttt{const}는 이 함수가 객체의 상태를 변경하지 않겠다는 약속이다.

\begin{lstlisting}
void showInfo() const
{
    cout << *number << endl;
}
\end{lstlisting}

읽기 전용 함수에 \texttt{const}를 붙이면 \texttt{const} 객체에서도 호출할 수 있다.

\begin{lstlisting}
const NumberBox box(10);
box.showInfo();
\end{lstlisting}

반면 값을 변경하는 함수에는 \texttt{const}를 붙이지 않는다.

\begin{lstlisting}
void setNumber(int value)
{
    *number = value;
}
\end{lstlisting}

\subsection{임시 포인터를 사용하는 이유}
복사 대입에서 다음 순서를 사용하면 새 메모리 할당이 성공한 뒤에만 기존 자원을 버리게 된다.

\begin{lstlisting}
int* newNumber = new int(*other.number);

delete number;
number = newNumber;
\end{lstlisting}

\texttt{newNumber}는 새 자원을 임시로 보관한다. 모든 준비가 끝난 뒤 \texttt{number}가 그 주소를 넘겨받는다. 이때 포인터 변수만 사라지고 실제 힙 메모리는 \texttt{number}가 계속 소유한다.

\subsection{전체 예제}
\begin{lstlisting}
#include <iostream>

using namespace std;

class NumberBox
{
private:
    int* number;

public:
    NumberBox(int value)
    {
        number = new int(value);
    }

    NumberBox(const NumberBox& other)
    {
        number = new int(*other.number);
    }

    NumberBox& operator=(const NumberBox& other)
    {
        if (this == &other)
        {
            return *this;
        }

        int* newNumber = new int(*other.number);

        delete number;
        number = newNumber;

        return *this;
    }

    ~NumberBox()
    {
        delete number;
    }

    void setNumber(int value)
    {
        *number = value;
    }

    int getNumber() const
    {
        return *number;
    }

    void showInfo() const
    {
        cout << "Number: " << *number << endl;
    }
};

int main()
{
    NumberBox box1(10);
    NumberBox box2 = box1;
    NumberBox box3(30);

    box2.setNumber(20);
    box3 = box1;
    box3.setNumber(40);

    box1.showInfo();
    box2.showInfo();
    box3.showInfo();

    return 0;
}
\end{lstlisting}

출력 결과는 다음과 같다.

\begin{verbatim}
Number: 10
Number: 20
Number: 40
\end{verbatim}

세 객체는 서로 다른 힙 메모리를 소유한다. 따라서 한 객체의 값을 바꾸어도 다른 객체의 값은 변하지 않는다.

\section{종합 정리}
\begin{importantbox}
\begin{itemize}[leftmargin=1.5em]
    \item 복사 생성자는 기존 객체로 새 객체를 만들 때 호출된다.
    \item 기본 복사는 멤버별 복사를 수행하며 포인터 멤버에서는 주소가 복사된다.
    \item 얕은 복사는 값 공유와 이중 해제 문제를 만들 수 있다.
    \item 깊은 복사는 새 힙 메모리를 할당하고 원본 값을 복사한다.
    \item 복사 대입 연산자는 이미 존재하는 객체 사이의 대입을 처리한다.
    \item 복사 대입에서는 자기 대입 검사와 기존 자원 정리가 필요하다.
    \item 소멸자, 복사 생성자, 복사 대입 연산자는 Rule of Three로 함께 설계한다.
    \item RAII는 객체의 수명과 자원의 수명을 연결한다.
    \item 자원 포인터는 \texttt{private}으로 숨기고 읽기 함수에는 \texttt{const}를 붙인다.
\end{itemize}
\end{importantbox}

\section{종합 실습}
\begin{practicebox}
\texttt{Book} 클래스를 작성한다.
\begin{itemize}[leftmargin=1.5em]
    \item \texttt{private} 멤버로 \texttt{int* price}를 가진다.
    \item 생성자에서 가격을 저장할 힙 메모리를 만든다.
    \item 깊은 복사를 수행하는 복사 생성자를 작성한다.
    \item 자기 대입을 처리하는 복사 대입 연산자를 작성한다.
    \item 소멸자에서 힙 메모리를 해제한다.
    \item \texttt{setPrice()}와 \texttt{showInfo() const}를 작성한다.
    \item 복사 생성과 복사 대입 후 객체들이 독립적으로 동작하는지 확인한다.
\end{itemize}
\end{practicebox}

\subsection{실습 완성 코드}
\begin{lstlisting}
#include <iostream>

using namespace std;

class Book
{
private:
    int* price;

public:
    Book(int price)
    {
        this->price = new int(price);
    }

    Book(const Book& other)
    {
        this->price = new int(*other.price);
    }

    Book& operator=(const Book& other)
    {
        if (this == &other)
        {
            return *this;
        }

        int* newPrice = new int(*other.price);

        delete this->price;
        this->price = newPrice;

        return *this;
    }

    ~Book()
    {
        delete this->price;
    }

    void setPrice(int price)
    {
        *this->price = price;
    }

    void showInfo() const
    {
        cout << "Price: " << *this->price << endl;
    }
};

int main()
{
    Book b1(30000);
    Book b2 = b1;
    Book b3(50000);

    b2.setPrice(40000);
    b3 = b1;
    b3.setPrice(100000);

    b1.showInfo();
    b2.showInfo();
    b3.showInfo();

    return 0;
}
\end{lstlisting}

\section{실습 파일 구성}
\begin{practicebox}
Day 7 실습 파일은 다음과 같이 관리한다.
\begin{itemize}[leftmargin=1.5em]
    \item \texttt{practice/p01\_copy\_constructor.cpp}
    \item \texttt{practice/p02\_shallow\_copy.cpp}
    \item \texttt{practice/p03\_deep\_copy.cpp}
    \item \texttt{practice/p04\_copy\_assignment.cpp}
    \item \texttt{practice/p05\_rule\_of\_three.cpp}
\end{itemize}
\end{practicebox}

\end{document}
